% Generated draft: requires mathematical and source review.
\documentclass[11pt]{article}
\usepackage[margin=25mm]{geometry}
\usepackage{fontspec}
\setmainfont{Latin Modern Roman}
\usepackage{amsmath,amssymb,mathtools,unicode-math}
\setmathfont{Latin Modern Math}
\usepackage{enumitem,needspace,xurl,hyperref}
\setlength{\parindent}{0pt}
\setlength{\parskip}{5pt}
\setlength{\emergencystretch}{3em}
\newcommand{\N}{\mathbb{N}}
\newcommand{\Z}{\mathbb{Z}}
\newcommand{\Q}{\mathbb{Q}}
\newcommand{\R}{\mathbb{R}}
\newcommand{\C}{\mathbb{C}}
\newcommand{\F}{\mathbb{F}}
\newcommand{\A}{\mathbb{A}}
\newcommand{\PP}{\mathbb{P}}
\newcommand{\E}{\mathbb{E}}
\newcommand{\eps}{\varepsilon}
\newcommand{\dd}{\,\mathrm{d}}
\begin{document}
\section*{UnsolvedMath ID: 20002010}
\textbf{Source problem code:} catalogue entry
\subsection{Definitions and mathematical statement}
Section C.5, ``Applications,'' contains two distinct questions.

(A) Arithmetic: for a nonzero polynomial \(f\in\Z[x]\) represented by a straight-line program of complexity \(\tau\), bound the number of roots congruent to \(1\) modulo a prime \(p\), as a function of \(\tau\). The aid reconstructs the intended quantity as
\[
 N_p(f)=\#\{a\in\Z:f(a)=0,\ a\equiv1\pmod p\},
\]
counting distinct integer roots, not residue classes, multiplicities or all \(p\)-adic roots. It interprets the intended target as a polynomial-in-\(\tau\) bound. For a standard division-free model one starts with \(1,x\) and uses additions, subtractions and multiplications; other constant-cost or division conventions change the input measure. The source does not specify whether \(p\) is fixed or the constants and exponent of a proposed polynomial bound must be uniform in \(p\). These versions are therefore not silently identified. It attributes amoeba-theoretic estimates to Maurice Rojas without stating an estimate or citation. No arithmetic bound is proved by the supplied partial result.

(B) Tropical geometry: formulate and prove a Calabi-type existence and uniqueness theorem for the tropical Monge--Amp\`ere measure in each specified metric class on a tropical Calabi--Yau. The source compares this with the real equation \(\det(\operatorname{Hess})=\text{constant}\), but defines neither the singular tropical space, metric class, operator, measure nor uniqueness normalization.

The aid treats a precise flat real-torus special case. Let \(\mathbb T^m=\R^m/\Lambda\), let \(A\) be a constant symmetric positive-definite matrix, and let \(u\in C^2(\R^m)\) be \(\Lambda\)-periodic with \(A+D^2u>0\). If
\[
 \det(A+D^2u)=\det A\quad\text{everywhere},
\]
then \(u\) is constant. Thus with, for example, \(\int_{\mathbb T^m}u=0\), the only solution is \(u=0\). The full potential is \(\tfrac12x^TAx+u(x)\), whose Hessian is periodic; the periodic correction alone cannot be globally strictly convex with positive constant Hessian determinant.

More generally, if \(u,v\) are periodic \(C^2\) potentials with \(A+D^2u>0\), \(A+D^2v>0\), and
\[
 \det(A+D^2u)=\det(A+D^2v)=\rho(x)>0,
\]
then \(u-v\) is constant. The aid uses matrix arithmetic--geometric mean and integration for constant density, and log-determinant linearization with the strong maximum principle for comparison. These are uniqueness/rigidity statements for existing smooth strictly convex potentials. They do not prove existence for general \(\rho\), define a singular tropical measure, or resolve the broad tropical Calabi conjecture.
\subsection{Short English statement}
The record asks for a straight-line-complexity bound on distinct integer roots in \(1+p\Z\), and for a tropical Calabi theorem. The aid clarifies the arithmetic interpretation and proves a flat real-torus uniqueness result with the necessary positive background Hessian. It supplies neither the arithmetic bound nor a general singular tropical existence theorem.
\subsection{Sources}
\begin{itemize}
\item Local UnsolvedMath registry, permanent ID \(20002010\), source code \texttt{catalogue entry}; canonical collection \texttt{aim-geometry-notes.json}, notes[347].
\item AIM workshop ``Amoebas and tropical geometry,'' Section C.5, ``Applications,'' parts A and B: \url{https://aimath.org/WWN/amoebas/amoebas.pdf}. No PDF was opened in this drafting session.
\item Maurice Rojas is credited for amoeba-theoretic arithmetic estimates, without an explicit bound or full authoritative reference. Original remarks and literature fields are empty; no proposer or proposed year is supplied.
\item Supplied research attempt \(1\), \texttt{partial\_result}, run status \texttt{partially\_solved}, rigor label ``proved,'' novelty confidence ``low.'' Registry classification \texttt{PARTIAL-PROGRESS}, updated 14 August 2026. Its candidate contribution is the integer-root interpretation and flat-background Monge--Amp\`ere rigidity/comparison package.
\end{itemize}
\end{document}
